\chapter{Harvard Thinking: 物語をどう語るか}

この章は、James Wood、Sam Marks、Lauren Groff、Nick White のあいだで交わされたパネル対話を、Samantha Liney-Perfoss の司会のもとに蒸留したものであり、このコレクションでは LazyingArt LLC によって編まれている。ここには黒板はなく、したがって回復すべき文字どおりの教室的導出もない。それでも、この講義には確かな背骨がある。まず書くことに厳しい基準を課し、ついで起源、人物とプロット、書き手自身の執着と身体、失敗と推敲、そして最後に、生命力をめぐる読者の試験へと進んでいく。私たちはその順序を緊密にたどることにしよう。議論が明快な関係を示唆するところでは、それを図式的に記すが、そうした式は講義を明晰にするためのものであって、黒板に書かれた法則ではないことを忘れてはならない。

\section{冒頭の基準: 代償を伴う書き方}

講義は、異様なほど厳しく始まる。人物やプロットや craft の用語からは入らない。そもそも、どのような書き方なら他人に影響を与えうるのかを問うところから始まるのである。答えは、感情的に効く文章は、それ自体が感情をこめて書かれていなければならない、というものだ。これは励ましの言葉ではない。基準なのである。

簡潔に言えば、冒頭の主張は
\begin{equation}
\text{読者への感情的効果} \Rightarrow \text{書き手による感情的投資}.
\end{equation}
となる。講義はただちにこれを強める。重要なのは、机に向かう時間が費やされたかどうかだけではない。効く文章は、時間以上のものを要求する。緊張、露出、危険である。
\begin{equation}
\text{効いている文章} \Rightarrow C > T,
\end{equation}
ここで \(T\) は単なる費やされた時間、\(C\) は仕事が要求するより広い代償である。これを量的な不等式として読むべきではない。生きた文章は、そこにいただけでは済まない、という言い方なのである。

同じ点は、``all cylinders firing'' という比喩で言い直される。何かがうまく働いているとき、書き手は自分の感情的・知的能力の縁まで押しやられているように感じる。図式的には、
\begin{equation}
\text{うまく働いている物語} \Rightarrow E + I,
\end{equation}
ここで \(E\) は動員された感情的能力、\(I\) は動員された知的能力である。これもまた図式的な記法にすぎない。講義が言っているのは、よい文章は純粋に感じられたものでも、純粋に考えられたものでもなく、その両方が同時に働いている状態を要する、ということだ。

この強い冒頭のあとではじめて、司会者は視野を広げる。物語ることは人間の生の一部である。ページと格闘する書き手という、なじみ深い像がある。もちろん craft もある。だが講義は、craft を全体と取り違えてほしくない。そこで、まさに正しいタイミングで、本当の問いが現れる。何が物語を忘れがたいものにするのか。

\subsection{Question \& Answer}

\paragraph{Question.} その文章が生きているかどうかについて、この講義が最初に課す試験とは何か。

\paragraph{Answer.} 最初の試験は、単なる polish ではない。書くことが、時間以上の何かをその書き手に代償として払わせているかどうかであり、感情的努力と知的努力の双方が、十分に動員されているかどうかである。講義は、ページ上の生気には、書くという行為そのものの生気が必要だと主張するところから始まる。

\section{物語はどこから始まるのか}

基準が置かれると、司会者は最初の実践的な問いを発する。物語はどこから始まるのか。パネルは、それに単一の起源を与えることを拒むことで答える。ある書き手にとっては、種は何年も抱えられているかもしれない。別の書き手にとっては、書き始めるまでは、物語は本当の意味では存在しない。さらに別の書き手にとっては、冒頭の一行や結びの一行が、中ほどがまだ見えていなくても、全体の構造を方向づけることがある。

とりわけ明晰な一つの説明では、長く温められてきたアイデアが、ある後日の出来事、読書、あるいは経験にぶつかるまでは不活性のままである、と語られる。その衝突を経てはじめて、それは物語の素材になるに足る切迫感を帯びる。これを要約すれば、
\begin{equation}
L + X \Rightarrow U + D + G + W,
\end{equation}
となる。ここで \(L\) は長く保持されたアイデア、\(X\) は後から起こる衝突、\(U,D,G,W\) は切迫感、密度、重力、重みを表す。講義の力点は、素材をただ持っている状態から、形を要求するほど強い圧力への移行に置かれている。

第二の説明は、矛盾というより補完である。物語は、あらかじめ蓄えられていて、ただ書き写されるのを待っているわけではない。それは書くという行為の中で、それ自身になっていく。
\begin{equation}
\text{熟成} \to \text{最初の執筆} \to \text{展開}.
\end{equation}
これは、この講義が、素朴な inspiration 観に加える最も重要な修正の一つである。物語は、あらかじめ完全に分かっているとは限らない。しばしば、それは書かれているあいだにしか可視化されない。

第三の説明は、方向づけを導入する。たとえ中ほどが不明でも、最初の一行、冒頭、あるいは結末が分かっていれば、それだけで空間を安定させるには十分である。
\begin{equation}
(\text{既知の冒頭}) \lor (\text{既知の結末}) \Rightarrow \text{物語的な方向づけ}.
\end{equation}
この点が重要なのは、まったく異なる書き手たちが、``どこから始めるかは分かっている'' と感じるとき、実際にはその言葉がかなり異なることを意味していても不思議ではない、ということを説明してくれるからである。

この部分の講義で最も実践的なエンジンは、人物から来る。ある書き手は、書くことはしばしばある特定の人物から始まり、そして二つの即座の問いから始まると言う。その人物は何を欲しているのか。そして何がそれを妨げているのか。
\begin{equation}
\text{物語のエンジン} = \text{人物} + \text{欲望} + \text{障害}.
\end{equation}
これは局所的ではあるが本当に有用な式である。すでに圧力が含まれており、そして圧力こそが、物語に運動を与えるものだからである。

\paragraph{A Worked Example.} パネルが与えるいくつかの起源の説明は、普遍的な algorithm を成すかのように偽ることなく、組み合わせることができる。
\begin{align}
L &\text{ は長いあいだ不活性なままでありうる},\\
L + X &\Rightarrow U + D + G + W,\\
U + D + G + W &\Rightarrow \text{最初の執筆},\\
\text{最初の執筆} &\Rightarrow \text{展開}.
\end{align}
これは storytelling の法則ではない。熟成、起動、最初の執筆、出現という、講義が共有していたパターンを注意深く再構成したものである。

この節は、重要な否定的教訓で閉じる。書き手たちは、恣意的なひねりで弱い草稿を救い出せるという幻想に警告を発する。``そして次にこんなとんでもないことが起こる'' と考えた瞬間、その結果はしばしば破綻である。物語の起源は、単なる装飾ではなく、獲得されたものであるように感じられなければならない。

\subsection{Question \& Answer}

\paragraph{Question.} 物語は実際にはどこから始まるのか。

\paragraph{Answer.} この講義の答えは複数形である。物語は長い熟成から始まることもあれば、古いアイデアと新しい経験との衝突から始まることもあり、最初の一行から始まることも、結末から始まることも、圧力のかかった人物から始まることもある。だが、いくつかの説明に共通している決定的な瞬間は、それらの素材が執筆へ移り、その中で展開し始めるときである。

\section{人物、行為、そしてプロットのスペクトラム}

講義はここで焦点を狭める。人物がエンジンとして導入された以上、次の問いは、人物と行為がどのような関係にあるかである。議論はまず、経験に近いところにとどまる。人物は書き手を驚かせることがあり、徐々に自らを明らかにすることもあり、ある場面ではよく見えても、別の場面ではそうでないこともある。劇作では、人物は見える以上に聞こえる存在かもしれない。二人の人物を同じ場面に置くと、それは実験の中の要素のように振る舞いうる。何が火花を散らすかを見るわけである。

この実験という言葉づかいが重要なのは、それがただちに理論的な極へつながるからである。Aristotle は、人物が行為に従属するという古典的見解の代表として召喚される。
\begin{equation}
\text{人物} \prec \text{行為}.
\end{equation}
パネルは、これを全面的な現代的プログラムとして採用するわけではない。むしろ、それを場の一端として用いるのである。現代の問題は、その原理に絶対的に従うことではなく、人物とプロットの双方が能動的に働いている物語空間をどう考えるかにある。

その空間は、スペクトラムとして記述される。
\begin{equation}
\text{人物主導} \leftrightarrow \text{プロット主導}.
\end{equation}
このスペクトラムの価値は診断的である。もしプロットがあまりに強引に支配すれば、私たちは虚偽の推進力、恣意的なひねり、そして ``プロットのためのプロット'' を得ることになる。逆に、人物が物語的圧力を欠いたまま支配すれば、読者は何も起こっていないと感じ始めるかもしれない。したがって二つの失敗の型は、
\begin{equation}
\text{不誠実な押し込み} \;\;\leftrightarrow\;\; \text{物語的圧力を欠いた人物}.
\end{equation}
と書くことができる。講義の中心的主張は、書き手はこの両方を避けねばならない、ということだ。

滑稽な宇宙船の例は、この点を鋭くしている。ときとして書き手は、物語がそれを獲得したからではなく、本当は何が続くべきなのかを見つけ出す困難から逃れたいがために、荒唐無稽な出来事を導入する。講義の応答は微妙である。宇宙船は、そもそも八ページ目に属していないのかもしれない。むしろ、それは、物語が本当はそこから始まるべき場所を示しているのかもしれない。これは美しい修正である。悪い逃避を、構造の問いへ変えるからだ。

この節はまた、人物が固定された内的な核ではないことも強調する。環境が重要なのである。同じ二人が同じ別れ話をするとしても、それがどこで起こるかによって場面は変わる。
\begin{align}
(\text{同じカップル},\text{浜辺}) &\Rightarrow \text{ひとつの行動の場},\\
(\text{同じカップル},\text{Everest}) &\Rightarrow \text{別の行動の場}.
\end{align}
したがって、人物とは、その人が ``何であるか'' だけではない。その人が、状況、場所、圧力、世界のもとで何になるかでもある。

そのあと講義は、印象的な促しを加える。自分が怖くてできないことを、人物にやらせてみよ。ここが転回点である。私たちは、物語の mechanics から、書き手と作品とのより深い関係へと移ってきたのだ。

\subsection{Question \& Answer}

\paragraph{Question.} 人物とプロットを、一方を他方へと平板化せずに、どのように考えるべきか。

\paragraph{Answer.} この講義は、それらを互いに必要としあう圧力の形として扱う。人物は欲望、障害、応答を与え、プロットはその圧力に方向と帰結を与える。失敗が生じるのは、プロットが恣意的になったとき、あるいは人物が物語的な力を失ったときである。よい storytelling は、その両方を同時に生かし続ける。

\section{魂、執着、場所、そして身体}

ここで司会者は、ふたたび問いを広げる。もし講義が人物主導の書き方とプロット主導の書き方について語ってきたのだとすれば、書き手の魂に導かれる書き方のようなものもありうるのだろうか。これはまさに正しいタイミングでの問いである。なぜなら、前節ですでに、書き手自身の恐れや勇気が人物の内部で働いていることが示唆されていたからだ。

答えは、執着、場所、そして繰り返し立ち返る内的な問いへ向かっていく。ある書き手は、初期の小説を、文字どおりの自伝ではないにせよ author-driven だったと語る。別の書き手は、自分が書く異なる作品は、より深い水準では、同じ内的問題の異なる版なのだと言う。場所は異様な力を伴って入ってくる。Mississippi のような地域は、外から選ばれた setting にすぎないのではない。骨の中に入り込んでいるものなのである。そこから逃れようとして何年も過ごし、やがてまたそこへ書き戻ろうとして何年も過ごすことになるかもしれない。

この節における最も力強い形式的修正は、
\begin{equation}
\text{虚構としての距離} \not\Rightarrow \text{自己の不在}.
\end{equation}
というものである。物語は、書き手の生きられた状況から遠く離れていても、なお深く個人的でありうる。これが重要なのは、講義が fiction を二つの反対方向の誤解から守ろうとしているからだ。個人的なものは文字どおりでなければならないという誤解と、創作されたものは非人格的でなければならないという誤解である。

遠い素材への橋は、単に biography 的なのではなく、身体的なものである。感覚記憶がアクセスを可能にする。
\begin{equation}
\text{感覚記憶} \Rightarrow \text{遠い虚構世界へのアクセス}.
\end{equation}
講義の具体例は見事である。コミューンに暮らしたことがなくても、手づくりの器に入った温かいオートミールを両手で包む感触なら知っている必要はない。その時点ですでに、それは身体的知識であり、そして身体的知識は転写できる。

ここでパネルの言葉づかいは、より比喩的になる。執着、自己の最も深い部分、書き手の魂、各人物を prism のように映す hologram、読者に与えられる近さの titration。私たちはそうした言葉を保存すべきだが、読み誤ってはならない。主張が神秘的な曖昧さにあるのではない。fiction は、繰り返し立ち現れる感情と記憶の構造に依拠しており、それらは新しい形、新しい時代、新しい setting へ翻訳されても生き残る、ということなのだ。

ここには重要な実践的帰結もある。書き手が何度も何度もある場所や、ある恐れや、ある問いに立ち返るとしても、それは想像力の失敗である必要はない。むしろ、作品が自らの本当の圧力場を見つけたしるしかもしれない。

\subsection{Question \& Answer}

\paragraph{Question.} fiction は、文字どおりの自伝でなくても、どうすれば深く個人的でありうるのか。

\paragraph{Answer.} この講義の答えは、自己は fiction の中へ、執着、場所、身体記憶、恐れ、欲望、そして繰り返し立ち返る内的問いを通して入ってくる、というものである。これらは、文字どおりの自己報告を必要としない。書き手は外的な状況を創作しながら、なお自分にとって最も深い素材から働きうるのである。

\section{失敗、草稿、推敲、そして読者}

講義はいま、process について最も細部に踏み込んだ節へ達する。失敗は、事後に隠されるべき恥としては扱われない。作業媒体そのものとして扱われる。ある書き手は、失敗を愛していると公言する。失敗そのものが快いからではなく、早すぎる清潔さなしに、遊び、やり直し、ほんとうの実験を可能にしてくれるからである。

したがって、講義の process 関係は、
\begin{equation}
\text{草稿} \to \text{失敗} \to \text{やり直し} \to \text{物語が自らを教える}.
\end{equation}
と書くことができる。その論理は精密である。草稿は不完全な素材を生む。不完全さは、ただ耐え忍ばれるのではなく、用いられる。やり直しは単なる繰り返しではなく、より明確な制約のもとでの別の試みである。十分な回数を経ると、物語は自分が何を要求しているかを、よりはっきりと示し始める。

ついでパネルは、作業の方法を比較する。ある書き手は notebooks を一度も読み返さず、ほとんど最後まで完全に analog のままでいる。別の書き手は、legal pad、打ち込まれたページ、printout、修正のあいだを行き来する。また別の書き手は、日々の生活を通して物語を持ち運ぶことを強調する。書くことは、机に向かって起こることだけではなく、歩きながら、悩みながら、自分で書いてしまった結び目をほどこうとしながら持続するものでもある。こうした方法は異なるが、講義はそれらを教条として競わせない。共通しているのは、process は意図的でなければならず、異なる process は素材との異なる接触のしかたを生む、という点である。

とりわけ力強い洞察は、revision が読者と結びつけられるところで現れる。熱のこもった草稿の最中には、まずい文章であってもすばらしく見えてしまうことがある。それはまだ書き手自身の切迫感に帯電しているからだ。revision は、そこにもう一つの心を導入する。だから、物語の beginning は、単に作品の時間的な最初の部分なのではない。
\begin{equation}
\text{冒頭} \Rightarrow \text{読者に、その物語の読み方を教える}.
\end{equation}
このために、冒頭は遅いのである。それは単なる始まりではない。指示なのである。

\paragraph{Derivation of the Revision Principle.} この講義の論理は、段階を追って次のように述べられる。
\begin{enumerate}
\item 初期の草稿はほとんど私的であり、その時点では、書き手だけが唯一の読者である。
\item revision は、たとえ想像された読者でしかなくとも、別の読者を導入する。
\item 冒頭の段落は、その読者がテクストのリズム、尺度、注意のモードを学ぶ最初の場所である。
\item したがって revision は、冒頭にとりわけ強く働きかける。
\item ひとたび冒頭が読者に正しく教えることができれば、その後の運動はより自由に進める。
\end{enumerate}

講義はまた、時間に対する sober さを強調する。草稿の熱気の中では、自分が何か見事なものを書いたと心から信じてしまうことがある。だが翌日、距離が回復すると、同じ箇所はまったく違って見えるかもしれない。それは inspiration への裏切りではない。判断が必要なかたちで到来しただけなのだ。

\subsection{Question \& Answer}

\paragraph{Question.} 失敗は、書く過程の中で何をしているのか。

\paragraph{Answer.} 失敗は、現在の試みの限界を明らかにする。それが作品の外側にある屈辱ではなく、作品の一部として受け入れられるとき、それは生産的になる。やり直しを促し、要求を明確にし、やがて物語がどんな形を望んでいるのかを、書き手自身に教え始めるのである。

\section{生命力、読むこと、編集、そして最後の助言}

最後の運動は、何が物語をよいものにするのかを問う。最初の答えは、生産的な意味で否定的である。よさは、まずその物語が realist か experimental か、連続的か断片的か、慣習的か奇矯かにあるのではない。まず vitality にある。
\begin{equation}
\text{よい物語} \Rightarrow \text{生命力} = \text{死んでいないこと} = \text{引き込まれること}.
\end{equation}
これが講義の終盤の支配的基準であり、冒頭の ``代償'' の基準に応答している。書き手に何かを支払わせた作品は、ページの上に生を生み出すはずなのである。

講義はいま、読むことへ向きを変える。読むこと自体が二段階に分けられる。
\begin{equation}
\text{honeymoon read} \to \text{seven-year itch}.
\end{equation}
第一の reading では判断を保留し、作品が自分に作用するに任せる。第二の rereading では分析的に読み返し、なぜある箇所が重要だったのかを問い、作品についての議論を組み立て始める。これは critic の方法としても、workshop の方法としても提示される。

書き手自身にとっての vitality の類比物は、身体的な基準である。作品が、それ自身の望むものに近づいてくるとき、人はそれを感じることができる。
\begin{equation}
\text{近づいてくる} \Rightarrow \text{それを感じられる}.
\end{equation}
講義はこれを vibration と呼ぶ。これは比喩的な言い方だが、空疎な言い方ではない。ある種の形式的成功は、明示的な理論によって知られる前に、まず接触によって知られる、という意味なのである。

まさにこの地点で、講義は冒頭の言葉へ戻る。感情的に効く文章は感情をこめて書かれねばならず、それは時間以上の代償を要求し、書き手を自分の能力の縁まで押しやる、とふたたび言われる。この回帰は重要である。円環を閉じるからだ。ページ上の vitality は、それを生み出す代償から切り離されてはいない。

次の実践的な試験は、音読である。音読は、テクストを私たちから疎外し、私たちをそれに対するより冷たい裁き手へ変える。
\begin{equation}
\text{音読} \Rightarrow \text{疎外} \Rightarrow \text{bullshit detector}.
\end{equation}
黙読の rereading では許せた文が、公の speech の中では即座に破綻することがある。講義の要点は、そのような失敗を言い逃れてはならない、ということである。それは証拠なのだ。

そこから自然に ``murder your darlings'' へつながる。だがここでもまた、講義は、判断が必要とされるところで slogan へ逃げることを拒む。
\begin{equation}
\text{うまく編集する} \neq \text{機械的に削る}.
\end{equation}
ある素材はたしかに自己満足であり、去らなければならない。だが別の素材は、まだ十分によく書かれていないから難しいだけなのかもしれない。だからこそ、この講義で最良の編集的な印は、``ここを削れ'' ではなく、``ここをもっとよくせよ'' なのだ。よい editing は、生命力の弱い箇所を見つけ出し、ただ原稿を短くするのではなく、より強い一文を要求する。

最後の助言は、すべてを習慣へ翻訳する。作品との接触を保ち続けよ。
\begin{equation}
\text{毎日作品に触れる} \Rightarrow \text{物語との接触が持続する}.
\end{equation}
writing とは何よりもまず活動であり、後から所有されるものではない。
\begin{equation}
\text{writing} = \text{process} \quad \text{before it is} \quad \text{product}.
\end{equation}
そして最後に、人は持ち込まれた rules によって硬直してはならない。講義は、要するに、最終的に重要なのはその物語自身が生み出す rules だけだ、と言って終わる。
\begin{equation}
\text{rules} = \text{いま手元にある物語に固有の rules}.
\end{equation}
これは怠惰のための許可証ではない。より大きな注意を要求するのである。物語自身の rules が聞こえるほどに、物語へ耳を澄まさなければならない。

\subsection{Question \& Answer}

\paragraph{Question.} 何が物語をうまく働かせるのか。そして、revision がそれを鈍らせるのではなく助けているとき、それはどう見分けられるのか。

\paragraph{Answer.} 物語がうまく働くのは、それに vitality があるとき、すなわち生きているように感じられ、私たちを引き込むときである。revision が助けになるのは、それがその生命を増すときであり、それが自己満足を取り除くことによってであれ、偽った音色をあらわにすることによってであれ、まだ自らを獲得していない箇所を強めることによってであれ、そうである。revision が失敗するのは、それが renewed attention ではなく、機械的な rule-following になってしまうときだけである。

\section{まとめ}

この講義は、代償から始まり、生命力で終わる。その二点のあいだで、storytelling についての首尾一貫した説明が展開される。物語は何年も熟成されることがあり、衝突によって生き始めることがあり、人物や冒頭の一行から始まることがあり、そして執筆という行為の中でしか本当に出現しないこともある。人物とプロットは均衡を要し、自己は執着、場所、身体記憶を通して fiction に入り、失敗は作品の外にあるのではなく、その媒体の一部であり、revision は読者を導入し、判断は rereading、音読、編集的距離を通して到来する。

講義全体を最後に一つの系列へ圧縮するなら、それは次のようになる。
\begin{equation}
\text{圧力} \to \text{執筆} \to \text{失敗} \to \text{revision} \to \text{生命力}.
\end{equation}
この連鎖がすべてではない。だが、それは議論全体を貫く力の線である。物語は圧力の中で始まり、言語の中で形をとり、失敗をくぐり抜け、revision から学び、そして最後に、ページの上で生きているかどうかによって試されるのである。